\section{Introduction}\label{sec:introduction}

Vision-language-action (VLA) driving models generate driving plans from visual observations and language inputs and explain their decisions in natural language~\cite{r1,r2,r3}. Alpamayo-R1 represents the causal grounds linking observations to a driving decision as Chain of Causation (CoC) reasoning and learns these explanations together with trajectory prediction~\cite{r1}. Such explanations help readers understand a model's judgment. Evaluating that judgment also requires checking whether the chosen behavior complies with the applicable rules.\par

This separate assessment is necessary because admissible behavior depends on the rules as well as the scene. Different waiting rules or stop-release conditions can permit different actions in the same observed situation. A driving model must therefore select an action appropriate to the scene and adjust that selection to the applicable rule. DriveReg addresses this requirement by retrieving situation- and region-specific regulations and assessing actions for legal compliance~\cite{r7}. Compliance is also a practical concern in deployed systems: in 2025, the U.S. National Highway Traffic Safety Administration opened an investigation following reports of red-light violations and wrong-way driving while Tesla's Full Self-Driving driver-assistance system was engaged~\cite{NHTSA2025PE25012}. Conditioning choices on applicable rules is thus one aspect of driving-model reliability.\par

For CoC-based models, this requirement raises the question of how to convey both the reason for a decision and the limits on admissible behavior. CoC explains why observations lead to a decision, but it need not express the applicability conditions and numerical parameters of an external rule as a case-specific constraint. The same scene and causal explanation can remain valid under different speed bounds, waiting durations, or release conditions, even though the permitted actions change. An action consistent with its explanation can therefore violate a separately supplied rule. Providing that rule at inference makes the relevant information available to the model, but its effect on behavior still needs to be assessed.\par

We assessed this effect in preliminary Alpamayo-R1-10B diagnostics without additional training. Explicit speed limits that the current speed already satisfied were still exceeded within two seconds in 6/9 PhysicalAI~\cite{nvidia2025physicalai} and 9/9 converted nuScenes~\cite{caesar2020nuscenes} open-loop outputs. Each source contributed three inputs evaluated with three trajectory-generation seeds. Raising the limit produced mean XY trajectory differences of at most 0.087 m over the full 6.4 s prediction horizon, averaged over three seeds per input and rule-input route. When asked to remain within 0.2 m for two seconds, all three predictions for one low-speed input from each source exceeded that radius. The violation counts were the same whether the rule was supplied in the user prompt or appended to CoC. In these small-sample diagnostics, explicitly stating the rule at inference was insufficient for compliance.\par

This observation motivates examining whether training with rule-dependent choices improves how a model uses supplied constraints. We evaluated this effect in a separate controlled synthetic candidate-selection study with Qwen3-VL-2B-Instruct~\cite{qwen2025qwen3vl2b}, distinct from the Alpamayo trajectory-generation diagnostics. Across three training seeds, with cross-entropy loss alone, identical initialization, and the same rule sentence at evaluation, CoC-only-trained adapters violated the rule in 17.92\% of pedestrian-entry predictions, compared with 1.11\% for adapters also trained with constraints. On pairs that fix the scene and candidates and change only the rule, the rate of answering both correctly rose from 1.94\% to 96.94\% (Section~\ref{sec:results-training}). These results show that constraint sentences can serve as learning information for rule-dependent selection in this task and support supplying them during training as well as inference.\par

Once constraints serve as learning information, errors in those constraints can also affect learning. Their correspondence with the case and intended rule therefore becomes a development concern. Constructing and conveying correct constraints presents three difficulties. First, a single constraint combines evidence with different roles, such as an observed clearance time, a supplied waiting duration, and an assumption that the zone remains clear. A sentence alone makes these sources and potential fault locations difficult to distinguish. Second, readable constraints can still be attached to the wrong cases or interpreted incorrectly in logic. Shuffling case--constraint associations during training left accuracy near 50\%, the constraint-free level (Section~\ref{sec:results-association}). Third, checking requires a formal representation, whereas the model learns from natural language. Specification checks must therefore connect to the generation of learning sentences, and the model's interpretation of those sentences must be assessed separately through behavior.\par

Existing research provides foundations on both sides of this connection. Traffic-rule formalization and verification and natural-language formalization give constraints explicit meanings that can be checked~\cite{r6,Wan2024,r8,r9}. Instruction tuning and natural-language principle learning improve models' use of language requirements~\cite{Wei2022,Bai2022,Guan2024}, and our preceding work improved action selection by appending natural-language guards to CoC~\cite{r4}. We connect these foundations in the development of case-specific learning constraints: their sources and meaning are checked, and the checked content is conveyed as natural-language input for learning.\par

We propose GuardSynth to implement this development path. GuardSynth first binds each constraint field to an observation, a supplied rule, or a modeling assumption. It then represents applicability, action restrictions, persistence and release, timing, and evidence in Evidence-Bound Lifecycle Contracts (EBLC). Supported contract profiles are lowered to Satisfiability Modulo Theories (SMT) queries that check consistency and designated properties within a restricted logical model. The checked specification is rendered as Controlled Natural Language (CNL) appended to the original CoC. Supported fixed-grammar parsers recover the fields from generated sentences to assess their correspondence with the specification. This path makes the relationships among sources, logical interpretations, and learning sentences explicit targets for checking.\par

The evaluation of GuardSynth distinguishes the behavioral effect of including constraints in training from the checkability of the constraints themselves. For the behavioral evaluation, we first compared fresh adapters with identical initialization and the same CNL at evaluation. Adapters trained with CNL using cross-entropy loss alone achieved violation rates of at most 1.56\% and rule-pair accuracy of at least 86.98\% on pedestrian entry and two position-hold representations. Adding semantic loss without rule text reduced violations but did not reach this level of rule-pair accuracy, nor did an 8B model given the same rule without task-specific training. Across six behavioral tasks, both natural-language constraints (P1) and CNL generated from EBLC (P2) improved rule-dependent selection. The behavioral gains came mainly from supplying constraints correctly associated with each case, and P2 did not outperform P1 under every condition.\par

Separately from these behavioral gains, we evaluated how well the EBLC path detects errors in constraints. In a fault-injection evaluation of the pedestrian-entry and position-hold specifications used to generate learning inputs, the checking system detected all 288 errors of six supported types and accepted all 48 valid cases. It did not detect the 48 cases in which the evidence record and specification were changed together to the same wrong value (Section~\ref{sec:results-checks}). Because source- and sentence-correspondence checks were added after training for this specification evaluation, we do not interpret the reported behavioral improvements as an effect of verification itself. These results distinguish errors that internal correspondence checks can detect from aspects requiring independent confirmation, such as observation validity and rule applicability.\par

Our contributions are:\par

\begin{enumerate}
\item A procedure that selects behavioral constraints relevant to a CoC learning case and binds each field to an observation, a supplied rule, or a modeling assumption.
\item An EBLC-based development path from source-bound specifications through restricted SMT checking and field--sentence correspondence checks to CNL appended to the original CoC.
\item An evaluation that separates the effect of training with constraints from the quality of the constraints themselves. It measures rule-dependent selection under identical initialization and evaluation inputs and examines case--constraint associations, wording changes, and the scope of checking for supported specification faults.
\end{enumerate}

Section~\ref{sec:related-work} reviews related work, and Section~\ref{sec:method} presents the problem setting and GuardSynth. Sections~\ref{sec:experimental-setup} and~\ref{sec:results} describe the experimental setup and results. Section~\ref{sec:discussion} discusses implications and limitations, and Section~\ref{sec:conclusion} concludes.
